\section*{Chapter \ref{ch:iv:research-case-studies-and-take-homes} --- Research Case Studies and Take-Homes}

\begin{solution}{iq:iv:research-case-studies-and-take-homes:1}
Row counts per stock and day (gaps, a short day); duplicated keys; bid above ask and zero or negative spreads; trade prices
outside the quotes; jumps in the daily close that match split ratios or dividends; time zone and session boundaries; the
distribution of each field for outliers. Count what each check finds, fix it or drop it, and record the decision.
\iqlookfor{a systematic, counted checklist before any analysis.}
\end{solution}

\begin{solution}{iq:iv:research-case-studies-and-take-homes:2}
First page: the question, the data checks and fixes with counts, the result with its standard error and one chart, the
robustness checks, the economic significance, the limits and next steps. Appendix: code that runs from a clean checkout,
further tables, the models that were tried and discarded (listed, so the reviewer can judge multiple testing).
\iqlookfor{a first page that stands alone and an appendix that makes the work reproducible.}
\end{solution}

\begin{solution}{iq:iv:research-case-studies-and-take-homes:3}
``Does the order-book imbalance at the end of a minute predict the mid-price return over the next minute, pooled across the
twenty stocks?'' It uses the data's richest field, has a clear economic mechanism (pressure on one side of the book), a
testable sign, and an obvious economic check (the size of the effect against the spread). A vaguer question (``fit a model to
everything'') cannot be finished honestly in four hours.
\iqlookfor{a precise, falsifiable question with a mechanism and a cost check.}
\end{solution}

\begin{solution}{iq:iv:research-case-studies-and-take-homes:4}
As a result: the question, the test, its power (the smallest effect it could have detected), the checks that were run, and the
conclusion that no effect larger than that bound exists in this sample. Then what you would try next and why. A clean null
with its power is more useful to a reviewer than a weak positive with no robustness.
\iqlookfor{reporting a null with its power, not apologising for it.}
\end{solution}

\begin{solution}{iq:iv:research-case-studies-and-take-homes:5}
40 duplicated (stock, day, minute) records: drop the duplicates. 15 rows with the bid above the ask: swap the fields (or drop
them if the cause is unclear). One stock whose price halves between two days: an unadjusted 2-for-1 split; multiply earlier
prices by one half, or later ones by two, consistently. Missing the split does not affect within-day minute returns, but any
daily return, rolling feature or volatility estimate across the split date shows a spurious 50\% fall.
\iqlookfor{finding each defect by a check, fixing it, and knowing which analyses it would corrupt.}
\end{solution}

\begin{solution}{iq:iv:research-case-studies-and-take-homes:6}
Returns within a stock-day share shocks and the regressor is persistent within the day, so the errors are not independent;
cluster by stock-day (or by day across stocks if there are common shocks). $0.48/0.022 \approx 22$: the effect is real in the
statistical sense, and on the chapter's data it matches the planted 0.5 within its error. Whether it matters is a separate
question (\cref{iq:iv:research-case-studies-and-take-homes:10}).
\iqlookfor{the clustering rationale and the separation of statistical from economic significance.}
\end{solution}

\begin{solution}{iq:iv:research-case-studies-and-take-homes:7}
Trades alternate at random between bid and ask, so a trade-price return contains the change in the side, which reverses: the
bid-ask bounce. With half-spread $h$ and mid-return volatility $\sigma$, the trade-price return is $r + h(s_{t} - s_{t-1})$
with independent random sides, so the lag-one covariance is $-h^2$ and the correlation $-h^2/(\sigma^2 + 2h^2) = -9/43 \approx
-0.21$. Mid prices have no bounce, so their autocorrelation is zero. A reversal ``signal'' built from trade prices would be
this artefact.
\iqlookfor{Roll's model and the formula, and the practical lesson.}
\end{solution}

\begin{solution}{iq:iv:research-case-studies-and-take-homes:8}
About 30 minutes of data checks, 30 of choosing the question, 90 of one test with its robustness checks, 45 of writing and 15
of re-running the code from scratch. Leave out: a second model family, hyperparameter searches, and any analysis you could not
finish and check. State what was left out in the memo.
\iqlookfor{a plan weighted towards data and writing, and discipline about scope.}
\end{solution}

\begin{solution}{iq:iv:research-case-studies-and-take-homes:9}
The second. Margin of the first: ``best of five, no out-of-sample test, no costs, what was the data check?''. Margin of the
second: ``clear question, honest test, small effect correctly sized; next, split in time and by stock''. Reviewers hire for
judgement, and the first memo shows none.
\iqlookfor{the reviewer's perspective: judgement over model count.}
\end{solution}

\begin{solution}{iq:iv:research-case-studies-and-take-homes:10}
No: even in the top decile the expected move is about $0.5 \times 0.9 = 0.45$ basis points, less than a tenth of the 6 basis
points a round trip crossing the spread costs. It can still be worth money as an input: to decide when to post passive orders
(post on the side the imbalance favours), to skew a market maker's quotes, or to time the execution of trades made for other
reasons; the value is then measured in execution cost saved, not in standalone P\&L.
\iqlookfor{comparing the effect with costs, and finding where a small effect is useful.}
\end{solution}

\begin{solution}{iq:iv:research-case-studies-and-take-homes:11}
Each null cell exceeds $|t| = 3.1$ with probability about 0.0019; over 60 independent cells, $1 - (1 - 0.0019)^{60} \approx
0.11$. Bonferroni at 5\% requires $p < 0.05/60$, that is $|t| \ge 3.34$, which 3.1 does not reach. There is no evidence of a
calendar effect.
\iqlookfor{the family-wise chance and the corrected threshold.}
\end{solution}

\begin{solution}{iq:iv:research-case-studies-and-take-homes:12}
Show the estimate by stock (a dot plot with standard errors) and by week or day (a rolling estimate), and the pooled estimate
with each stock left out in turn. If the effect is spread across stocks and time, say so with the chart. If the reviewer is
right and one stock or one week drives it, concede, report the estimate without it, and reframe the conclusion as a finding
about that stock or episode, with the smaller sample's uncertainty.
\iqlookfor{robustness shown with data, and a gracious, quantified concession.}
\end{solution}
